\documentclass[hyperref={pdftex,unicode}]{beamer}

\usepackage[T2A]{fontenc}
\usepackage[utf8]{inputenc}
\usepackage[russian]{babel}
\usepackage{cmap}

\usepackage{xcolor}

% \usepackage{helvet}
\usepackage{pscyr}

\usepackage{multicol}

% \usepackage{amssymb,amsfonts,amsmath,mathtext}
% \usepackage{cite,enumerate,float}

\usepackage{listings}

\graphicspath{{images/}}
\usetheme{boxes}
\usecolortheme{whale}
\useinnertheme{default}
\useoutertheme{default}
\usefonttheme{professionalfonts}

\definecolor{Blue}{RGB}{55,110,160}
\definecolor{Yellow}{RGB}{255,210,65}
\definecolor{Green}{RGB}{75,120,60}
\definecolor{LightGray}{RGB}{205,205,205}

\setbeamercolor*{palette primary}{fg=white,bg=Blue}

\setbeamercolor*{enumerate item}{fg=Yellow}
\setbeamercolor*{enumerate subitem}{fg=Yellow}
\setbeamercolor*{enumerate subsubitem}{fg=Yellow}

\setbeamercolor*{itemize item}{fg=Yellow}
\setbeamercolor*{itemize subitem}{fg=Yellow}
\setbeamercolor*{itemize subsubitem}{fg=Yellow}

\lstset{
  language=python,
  basicstyle=\small\ttfamily,
  keywordstyle=\bfseries\color{Blue},
  commentstyle=\color{Green},
  numberstyle=\footnotesize\color{LightGray},
  numbersep=0.8em,
  stepnumber=1,
  keepspaces=true,
  breaklines=true,
  aboveskip=0.5\baselineskip,
  belowskip=0.5\baselineskip}


\title{Python: functions}
\author{meequz@gmail.com \\ budnyjj@pirates.by}
\date{}

\begin{document}

\begin{frame}
  \maketitle
\end{frame}

\begin{frame}{Сегодня}
  \begin{itemize}
    \item ООП
      \begin{itemize}
      \item Инкапсуляция
      \item Наследование
      \item Полиморфизм
      \end{itemize}
    \item Работа с исключениями
  \end{itemize}
\end{frame}

\begin{frame}{Объектно-ориентированное программирование}
  Объект --- сущность, состоящая из \\
  \begin{itemize}
  \item данных (полей, атрибутов)
  \item методов обработки этих данных (методов объекта).
  \end{itemize}
\end{frame}

\begin{frame}{Пример объекта}
  Кроме прочих своих достоинств, кот
  \begin{itemize}
    \item демонстрирует характерное поведение,
    \item реагирует на сообщения,
    \item наделён унаследованными реакциями,
    \item управляет своим, вполне независимым, внутренним состоянием.\footnote[frame]{
   Roger King, My cat is object-oriented}
 \end{itemize}
\end{frame}

\begin{frame}{Пример объекта}
 \begin{figure}[h!]
   \includegraphics[width=0.6\linewidth]{cat.jpg}
 \end{figure}
\end{frame}

\begin{frame}{Область применения ООП}
  Объектно-ориентированный подход хорош там, \\
  где проект подразумевает развитие.
\end{frame}


\begin{frame}[fragile]{Объявление класса [dataClass.py]}
    \begin{lstlisting}[numbers=left]
class exampleDataClass:
    '''A data example class'''

    cls_var = 'cls_data'

    def __init__(self):
        self.obj_var = 'obj_data'

    def f(self):
        return self.obj_var
    \end{lstlisting}
\end{frame}

\begin{frame}[fragile]{Атрибуты объекта [dataExample1.py]}
\begin{minipage}{0.62\linewidth}
    \begin{lstlisting}[numbers=left,basicstyle=\scriptsize\ttfamily]
import dataClass as dc

o1 = dc.exampleDataClass();
o2 = dc.exampleDataClass();

print('obj_var is object data member:')
print('BEFORE:')
print('o1.obj_var: ', o1.obj_var)
print('o2.obj_var:', o2.obj_var)

o1.obj_var = 'new_object_data'

print('AFTER:')
print('o1.obj_var:', o1.obj_var)
print('o2.obj_var:', o2.obj_var)
    \end{lstlisting}
\end{minipage}
\hfill
\begin{minipage}{0.35\linewidth}
    \begin{lstlisting}[numbers=left,basicstyle=\tiny\ttfamily,numbers=none]
class exampleDataClass:
    cls_var = 'cls_data'

    def __init__(self):
        self.obj_var =
            'obj_data'

    def f(self):
        return self.obj_var
    \end{lstlisting}
\end{minipage}
\end{frame}

\begin{frame}[fragile]{Атрибуты класса [dataExample2.py]}
\begin{minipage}{0.62\linewidth}
    \begin{lstlisting}[numbers=left,basicstyle=\scriptsize\ttfamily]
import dataClass as dc

o1 = dc.exampleDataClass();
o2 = dc.exampleDataClass();

print('Data is object data member:')
print('BEFORE:')
print('o1.cls_var:', o1.cls_var)
print('o2.cls_var:', o2.cls_var)

dc.exampleDataClass.cls_var = \
               'new_object_data'

print('AFTER:')
print('o1.cls_var: ', o1.cls_var)
print('o2.cls_var:', o2.cls_var)
    \end{lstlisting}
\end{minipage}
\hfill
\begin{minipage}{0.35\linewidth}
    \begin{lstlisting}[numbers=left,basicstyle=\tiny\ttfamily,numbers=none]
class exampleDataClass:
    cls_var = 'cls_data'

    def __init__(self):
        self.obj_var =
            'obj_data'

    def f(self):
        return self.obj_var
    \end{lstlisting}
\end{minipage}
\end{frame}

\begin{frame}[fragile]{Изменение атрибутов объекта}
\begin{center}
  \begin{tabular}[h!]{p{0.45\linewidth} p{0.54\linewidth}}
    getattr(object, name) & var = object.name \\
    setattr(object, name, value) & object.name = value \\
    delattr(object, name) & del object.name \\
    hasattr(object, name) & \scriptsize{
      implemented by calling getattr(object, name)

      and seeing whether it raises an exception or not} \\
  \end{tabular}
\end{center}
\end{frame}

\begin{frame}[fragile]{Изменение атрибутов объекта [dataExample3.py]}
\begin{minipage}{0.62\linewidth}
    \begin{lstlisting}[numbers=left,basicstyle=\scriptsize\ttfamily]
import dataClass as dc

o = dc.exampleDataClass()
print(dir(o))

del o.obj_var
o.new_data = 'some_data'
o.f2 = lambda x: x * 3

print(dir(o))
print(o.new_data)
print(o.f2('M'))
    \end{lstlisting}
\end{minipage}
\hfill
\begin{minipage}{0.35\linewidth}
    \begin{lstlisting}[numbers=left,basicstyle=\tiny\ttfamily,numbers=none]
class exampleDataClass:
    cls_var = 'cls_data'

    def __init__(self):
        self.obj_var =
            'obj_data'

    def f(self):
        return self.obj_var
    \end{lstlisting}
\end{minipage}
\end{frame}

\begin{frame}[fragile]{Инкапсуляция}
    \begin{lstlisting}
class Simple:
    '''Simple class with private attribute'''
    def __init__(self):
        self.__private_attr = 20
        print(self.__private_attr)

s = Simple()
print(s.__private_attr)
\end{lstlisting}
\end{frame}

\begin{frame}[fragile]{Наследование}
Синтаксис:
    \begin{lstlisting}[numbers=none]
class Derived(Base):
    pass
class Derived(module_name.Base):
    pass
class Derived(Base1, Base2, Base3):
    pass
    \end{lstlisting}
Особенности:
\begin{itemize}
\item Все методы --- виртуальные
\item Вызов метода базового класса: \lstinline$Base.method()$
\item Порядок поиска атрибута:
  \begin{enumerate}
    \item
      Derived
    \item Base1, затем рекурсивно в базовых классах
    \item Base2, затем рекурсивно в базовых классах $\dots$
  \end{enumerate}
\end{itemize}
\end{frame}

\begin{frame}[fragile]{Наследование}
\begin{itemize}
\item type(object)
\item isinstance(object, classinfo)
\item issubclass(class, classinfo)
\end{itemize}
\end{frame}


\begin{frame}{Встроенные методы классов}
  \begin{itemize}
  \item \textbf{\_\_doc\_\_}
  \item \_\_name\_\_
  \item \_\_module\_\_
  \item \_\_bases\_\_
  \end{itemize}
\end{frame}

\begin{frame}{Встроенные методы объектов}
  \begin{itemize}
  \item \textbf{\_\_init\_\_}
  \item \textbf{\_\_dict\_\_}
  \item \_\_class\_\_
  \item \_\_del\_\_
  \item \_\_cmp\_\_
  \item \_\_hash\_\_
  \item \_\_getattr\_\_
  \item \_\_setattr\_\_
  \item \_\_delattr\_\_
  \item \_\_call\_\_
  \end{itemize}
\end{frame}

\begin{frame}[fragile]{Эмуляция последовательностей}
  \begin{minipage}{0.3\linewidth}
    \begin{itemize}
    \item \_\_len\_\_
    \item \_\_getitem\_\_
    \item \_\_setitem\_\_
    \item \_\_delitem\_\_
    \item \_\_getslice\_\_
    \item \_\_setslice\_\_
    \item \_\_delslice\_\_
    \item \_\_contains\_\_
    \end{itemize}
  \end{minipage}
  \hfill
  \begin{minipage}{0.65\linewidth}
    \begin{lstlisting}[basicstyle=\scriptsize\ttfamily]
import logging

class LoggingDict(dict):
    def __setitem__(self, key, value):
        logging.info(
            'Setting {k} to {v}'.\
            format(k=key, v=value))
        return super(LoggingDict, self).\
        __setitem__(key, value)

logging.basicConfig(level=logging.INFO)

ld = LoggingDict()
ld['a'] = 123
      \end{lstlisting}
  \end{minipage}
\end{frame}

\begin{frame}{Приведение к базовым типам}
  \begin{itemize}
  \item \textbf{\_\_repr\_\_}
  \item \textbf{\_\_str\_\_}
  \item \_\_oct\_\_
  \item \_\_hex\_\_
  \item \_\_complex\_\_
  \item \_\_int\_\_
  \item \_\_long\_\_
  \item \_\_float\_\_
  \end{itemize}\
\end{frame}

\begin{frame}[fragile]{Пример перегрузки оператора}
    \begin{lstlisting}[basicstyle=\tiny\ttfamily]
class SignableMatrix:
    def __init__(self, array=[[]]):
        self.array = dc(array)
        dim_size = len(self.array)
        self.__row_signs = [False for i in range(dim_size)]
        self.__col_signs = [False for i in range(dim_size)]

    def __repr__(self):
        repr_str = 'Matrix: [\n'
        for row in self.array:
            repr_str += '    {}\n'.format(row)
        repr_str += ']\n'
        repr_str += 'Row signs:\n    {}\n'.format(self.__row_signs)
        repr_str += 'Column signs:\n    {}'.format(self.__col_signs)

        return repr_str

init_graph = [[1, 2, 3], [1, 2, 3], [1, 2, 3]]
init_matrix = SignableMatrix(init_graph)

print(init_matrix)
    \end{lstlisting}
\end{frame}

\begin{frame}{Остальные магические методы}
\begin{itemize}
  \item rafekettler.com/magicmethods.html
  \item \scriptsize{docs.python.org/2/reference/datamodel.html\#special-method-names}
\end{itemize}
\end{frame}

% \begin{frame}{Декораторы классов}
%   \begin{itemize}
%   \item @staticmethod
%   \item @classmethod
%   \item @property
%   \item @setter
%   \item @deleter
%   \end{itemize}
% \end{frame}

\begin{frame}{Исключения}
Обработка исключительных ситуаций --- \\
механизм языков программирования, \\
предназначенный для описания реакции программы на возможные проблемы (исключения), \\
которые могут возникнуть при выполнении программы \\
и приводят к невозможности (бессмысленности) \\
дальнейшей отработки программой её базового алгоритма.~\footnote[frame]{
http://en.wikipedia.org/wiki/Exception\_handling}
\end{frame}

\begin{frame}[fragile]{Пример: открытие файла}
  \begin{lstlisting}
for arg in sys.argv[1:]:
  try:
      f = open(arg, 'r')
  except IOError:
      print('cannot open', arg)
  else:
      print(arg, 'has',
            len(f.readlines()), 'lines')
      f.close()
  \end{lstlisting}
\end{frame}

\begin{frame}[fragile]{Пример: безопасное деление}
  \begin{lstlisting}[numbers=left]
def divide(x, y):
     try:
         result = x / y
     except ZeroDivisionError:
         print('division by zero!')
     else:
         print('result is', result)
     finally:
         print('executing finally clause')
  \end{lstlisting}
\end{frame}

\begin{frame}[fragile]{Как сгенерировать исключение}
  \begin{lstlisting}[numbers=left]
try:
    raise Exception('spam', 'eggs')
except Exception as inst:
    print(type(inst))
    print(inst.args)
    print(inst)
    x, y = inst.args
    print('x =', x)
    print('y =', y)
  \end{lstlisting}
\end{frame}

\begin{frame}[fragile]{Стандартные исключения}
  \begin{itemize}
  \item Exception
  \item ImportError
  \item IOError
  \item AttributeError
  \item ValueError
  \item SyntaxError
  \item ZeroDivisionError
  \item \dots
  \end{itemize}
\end{frame}

\begin{frame}{Полезные ссылки}
  \begin{itemize}
    \item docs.python.org/3/tutorial/classes.html
    \item docs.python.org/3/library/exceptions.html
    \item rafekettler.com/magicmethods.html
      \scriptsize{
      \item docs.python.org/3/reference/datamodel.html\#special-method-names
      \item www.ibm.com/developerworks/ru/library/l-python\_part\_6/
      \item www.ibm.com/developerworks/ru/library/l-python\_part\_7/
      }
  \end{itemize}
\end{frame}

\begin{frame}{Спасибо за внимание!}
  \begin{itemize}
  \item https://github.com/budnyjj/courses\_python
  \item https://vk.com/budnyjj
  \item budnyjj@gmail.com
  \end{itemize}
\end{frame}

\end{document}